\documentclass[12pt]{article}
\usepackage[margin=0.72in]{geometry}
\usepackage{fontspec}
\setmainfont{Latin Modern Roman}
\usepackage{microtype}
\usepackage{fancyhdr}
\usepackage{titlesec}
\usepackage{enumitem}
\usepackage{needspace}
\usepackage{xcolor}
\usepackage[hidelinks]{hyperref}
\setlength{\parindent}{0pt}
\setlength{\parskip}{0.34em}
\setlength{\headheight}{15pt}
\titleformat{\section}{\large\bfseries\scshape}{}{0pt}{}
\titleformat{\subsection}{\normalsize\bfseries}{}{0pt}{}
\pagestyle{fancy}
\fancyhf{}
\fancyhead[L]{Whately Logic: Week 3 Test}
\fancyhead[R]{The Professor's Riddles}
\fancyfoot[C]{\thepage}
\renewcommand{\headrulewidth}{0.4pt}

\newcommand{\focusbox}[1]{%
\par\vspace{0.35em}\noindent\begingroup\setlength{\fboxsep}{6pt}\colorbox{gray!12}{\begin{minipage}{\dimexpr\linewidth-2\fboxsep\relax}#1\end{minipage}}\endgroup\par\vspace{0.45em}}
\newcommand{\studentline}{\par\vspace{0.60em}\hrule\vspace{0.64em}}
\newcommand{\shortline}{\par\vspace{0.38em}\hrule\vspace{0.48em}}

\begin{document}
\begin{center}
{\LARGE\bfseries Week 3 Logic Test}\par\vspace{0.16em}
{\large\scshape The Professor's Riddles}\par\vspace{0.25em}
{Department of Science — Spotting Errors with AEIO Quantifiers}\par\vspace{0.45em}
\rule{0.82\textwidth}{0.4pt}
\end{center}

\section*{Name: \rule{2.35in}{0.4pt} \hfill Date: \rule{1.25in}{0.4pt}}

\focusbox{\textbf{Whately's guidance:} Before judging whether a claim is true, state its exact quantity and quality. Use the four standard forms: A = All S are P (universal affirmative), E = No S are P (universal negative), I = Some S are P (particular affirmative), O = Some S are not P (particular negative). "Most" is a strong particular claim and must be justified by evidence; it is not the same as "All." Overreaching with All/None when only Some holds, or claiming Most when evidence is unclear, produces false propositions even if the subject and predicate are right. Logic tests the structure of what is asserted.}

\section*{The Puzzle Story}

At the Academy of Straight Thinking, the Department of Science occupies the east wing. The professor has prepared four "riddles"—lists of statements about important scientific categories—for a classification activity. Each list is meant to contain only statements that could belong in a *true* description of that category.

The professor writes on the board:

``Tell me where my riddles go wrong. Some of these statements overreach. Use AEIO logic to find the errors in scope.''

Below are the four riddles. For each category, decide which statements could belong in a true riddle (they are accurate in their claims about quantity) and which are false because they use All, None, or Most too broadly or incorrectly.

\section*{Part I: The Four Riddles}

For each riddle, mark statements as T (could belong in a true riddle) or F (false due to overreach). Then, in the spaces provided after each riddle, list the false statement numbers and briefly explain the AEIO error (e.g., "claims All (A) but only Some (I) is true").

\subsection*{Riddle 1: Mammals}

Decide which statements are true and which are false.

\begin{enumerate}[leftmargin=*, itemsep=0.1em]
\item All are vertebrates.
\item Some lay eggs.
\item None can live in water.
\item Most have hair or fur at some point in life.
\item All give birth to live young.
\item Some can glide through the air.
\item None breathe through gills as adults.
\item Some are larger than cars.
\item All are warm-blooded.
\item Some use echolocation.
\end{enumerate}

\textbf{False statement numbers:} \hrulefill

\textbf{Explain the scope error(s) using AEIO terms for one or two false statements:}\studentline

\Needspace{0.25\textheight}

\subsection*{Riddle 2: Birds}

Decide which statements are true and which are false.

\begin{enumerate}[leftmargin=*, itemsep=0.1em]
\item All have feathers.
\item Some cannot fly.
\item None can swim.
\item All lay eggs.
\item Most can fly.
\item None have teeth as adults.
\item Some hunt at night.
\item All have hollow bones.
\item Some migrate thousands of miles.
\item None are warm-blooded.
\end{enumerate}

\textbf{False statement numbers:} \hrulefill

\textbf{Explain the scope error(s) using AEIO terms:}\studentline

\Needspace{0.25\textheight}

\subsection*{Riddle 3: Insects}

Decide which statements are true and which are false.

\begin{enumerate}[leftmargin=*, itemsep=0.1em]
\item All are invertebrates.
\item All have six legs as adults.
\item Some have wings.
\item None live in water.
\item Most are smaller than a human hand.
\item All undergo metamorphosis.
\item Some live in colonies.
\item None have antennae.
\item Some can sting.
\item All have exoskeletons.
\end{enumerate}

\textbf{False statement numbers:} \hrulefill

\textbf{Explain the scope error(s) using AEIO terms:}\studentline

\Needspace{0.25\textheight}

\subsection*{Riddle 4: Planets}

Decide which statements are true and which are false.

\begin{enumerate}[leftmargin=*, itemsep=0.1em]
\item All orbit a star.
\item None produce their own light.
\item Some have rings.
\item All are made of rock.
\item Some have no moons.
\item Most spin on an axis.
\item None have atmospheres.
\item Some are larger than Earth.
\item All are round because of gravity.
\item None can support storms.
\end{enumerate}

\textbf{False statement numbers:} \hrulefill

\textbf{Explain the scope error(s) using AEIO terms:}\studentline

\clearpage

\section*{Part II: Spiral from Week 1}

Week 1 emphasized asking what follows from what and distinguishing claims that can be tested. Choose one statement from any of the riddles above. Translate it into a precise AEIO proposition (name the form: A, E, I, or O). Then state what would follow if we accepted it as true.

\textbf{Statement chosen (riddle and number):} \hrulefill

\textbf{Precise AEIO form:}\shortline

\textbf{Subject:}\shortline

\textbf{Predicate:}\shortline

\textbf{If we accept this precise proposition, what (if anything) logically follows about the category? Briefly explain.}\studentline

\Needspace{0.35\textheight}

\section*{Part III: Spiral from Week 2}

Week 2 focused on fixing terms that shift meaning. Scan all four riddles. Are there any words that could be "two-faced" (carry different meanings in different statements or contexts)? If yes, identify one and fix it. If no clear cases, explain why the statements here keep terms steadier than the Week 2 examples.

\textbf{Two-faced word found (or "none") :}\shortline

\textbf{Fixed meaning for this test:}\shortline

\textbf{Why fixing (or the absence of shifting) matters before judging scope:}\studentline

\Needspace{0.35\textheight}

\section*{Part IV: Repair One False Claim}

Choose one false statement from the riddles (note which). Rewrite it as a true proposition using the correct AEIO quantifier. Then add one brief real-world observation that would support the repaired claim.

\textbf{Riddle and statement number:}\hrulefill

\textbf{Repaired true version (with correct quantifier):}\studentline

\textbf{One supporting observation or fact:}\studentline

\textbf{Why does the precise AEIO version make the claim easier to test or defend?}\studentline

\Needspace{0.30\textheight}

\section*{Part V: What Logic Can Decide Here}

For each item, mark whether \textbf{Logic (using AEIO) can help test or expose the error} or \textbf{Logic alone cannot settle this (needs facts outside logic)}.

\begin{enumerate}[leftmargin=*]
\item Whether "All mammals give birth to live young" overclaims because platypuses lay eggs. \hfill \rule{1.6in}{0.4pt}
\item Whether the subject term "planets" is distributed in a universal claim. \hfill \rule{1.6in}{0.4pt}
\item Whether insects in fact live in water as adults or larvae. \hfill \rule{1.6in}{0.4pt}
\item Whether changing "All are warm-blooded" to "Some are warm-blooded" changes the quality of the proposition. \hfill \rule{1.6in}{0.4pt}
\item Whether birds are the only animals with feathers (definition or fact). \hfill \rule{1.6in}{0.4pt}
\item Whether "Most" claims can ever be strict A or E propositions. \hfill \rule{1.6in}{0.4pt}
\end{enumerate}

\textbf{Explain your marking for two of the items above:}\studentline\studentline

\Needspace{0.30\textheight}

\section*{Final Whately Reflection}

In one or two sentences, explain why a scientist (or any thinker) must pay close attention to the exact words "all," "some," "none," and "most" before using a set of statements as a reliable riddle or description of a category.

\studentline\studentline

\end{document}